\newpage
\section{JORNADAS DE USUÁRIO}
A partir das personas geradas foram escritas as seguintes jornadas de usuário dentro de nossa plataforma. Com elas buscamos cobrir todas as principais funções do MVP.

\subsection{Jornada 1 - Maria (Sênior): Inscrição em aula ao vivo}
Objetivo: se inscrever em uma aula de ginástica que ocorrerá na próxima semana.
\begin{table}[H]
\centering
\caption{Jornada de usuário 1}
\label{tab:jornada1}
\begin{tabular}{|m{3.2cm}|m{5cm}|m{2.5cm}|m{4.5cm}|}
\hline
\multicolumn{1}{|c|}{\textbf{Etapa}} & \multicolumn{1}{c|}{\textbf{Ação}} & \multicolumn{1}{c|}{\textbf{Emoção}} & \multicolumn{1}{c|}{\textbf{Ponto de fricção / ganho}} \\ \hline

1 - Recebe lembrete & Recebe notificação avisando que uma nova aula de ginástica foi agendada para a próxima semana & Curiosa, insegura & + Notificação clara com nome da atividade, dia e horário \\ \hline

2 - Abrir o app login & Abre o VivaMais e realiza o login simplificado & Confiante & + Login simples \\ \hline

3 - Encontra a atividade & Visualiza no calendário da tela inicial o card da aula disponível & Levemente Insegura & + Card com data, horário e vagas disponíveis visíveis \newline\newline - Pode ter medo de não funcionar ou das implicações de inscrever-se. \\ \hline

4 - Realiza a inscrição & Toca em ``Inscrever-se'', então aparece uma tela de confirmação simples com os dados da aula e botão ``Confirmar'' & Tranquila & + Tela de confirmação evita inscrição acidental \newline\newline + Linguagem simples e direta na confirmação \\ \hline

5 - Confirmação recebida & Após confirmar, uma mensagem aparece: ``Inscrição confirmada!'' O evento fica salvo na aba de Minhas Atividades & Satisfeita & + Feedback imediato e positivo na tela \newline\newline + Atividade adicionada automaticamente a Minhas Atividades \\ \hline

6 - Recebe lembrete da aula & Chegando perto da sessão, recebe notificações lembrando do evento & Segura e preparada & + Lembretes automáticos \\ \hline
\end{tabular}
\end{table}

\subsection{Jornada 2 - Maria (Sênior): Participação de aula ao vivo}
Objetivo: Acessar aula de Yoga, na qual ela já está inscrita, e que começará em breve.
\begin{table}[H]
\centering
\caption{Jornada de usuário 2}
\label{tab:jornada2}
\begin{tabular}{|m{3.2cm}|m{5cm}|m{2.5cm}|m{4.5cm}|}
\hline
\multicolumn{1}{|c|}{\textbf{Etapa}} & \multicolumn{1}{c|}{\textbf{Ação}} & \multicolumn{1}{c|}{\textbf{Emoção}} & \multicolumn{1}{c|}{\textbf{Ponto de fricção / ganho}} \\ \hline

1 - Notificação & Recebe notificação no celular avisando que uma aula começará em breve & Curiosa, insegura & + Notificação em linguagem simples e amigável \newline\newline - Precisa encontrar um lugar para abrir a aula \\ \hline

2 - Abre o app & Encontra um lugar adequado e abre o aplicativo & Aliviada & + Acesso simples, em até 2 toques \\ \hline

3 - Faz login & Realiza o login simplificado na plataforma, no caso do app no celular poderia já estar conectada e abrir direto & Confiante & + Login simplificado \newline\newline - Pode ter problemas com a senha \\ \hline

4 - Entrar na aula & Vai para aba ``Ao Vivo'' e entra na reunião ou na aba ``Atividades'' vê o card de yoga com etiqueta ``Entrar Agora'' em destaque & Levemente Insegura & - Pode ser um pouco difícil se localizar nas primeiras tentativas \newline\newline + Vários pontos de acesso \\ \hline

5 - Participar & Entra na reunião, escuta a tutora e os colegas e pode participar & Interessada & - Pode haver problemas com câmera ou microfone \\ \hline

6 - Finalizar & Participa de toda a sessão e sai quando acabada & Satisfeita & + Volta para tela de início da plataforma \\ \hline
\end{tabular}
\end{table}

\subsection{Jornada 3 - Maria (Sênior): Disponibilização de aula ao vivo}
Objetivo: agendar uma nova roda de conversa e acompanhar quem participou.
\begin{table}[H]
\centering
\caption{Jornada de usuário 3}
\label{tab:jornada3}
\begin{tabular}{|m{3.2cm}|m{5cm}|m{2.5cm}|m{4.5cm}|}
\hline
\multicolumn{1}{|c|}{\textbf{Etapa}} & \multicolumn{1}{c|}{\textbf{Ação}} & \multicolumn{1}{c|}{\textbf{Emoção}} & \multicolumn{1}{c|}{\textbf{Ponto de fricção / ganho}} \\ \hline

1 - Acesso e login & Entra na plataforma e realiza o login utilizando suas credenciais de tutora & Confiante & + Login simples, mesmo para tutores \\ \hline

2 - Cria sessão & Clica em ``Adicionar nova sessão'', preenche nome, data, horário, duração e limite de participantes & Produtiva & + Formulário simples \\ \hline

3 - Notificação automática & Sistema envia automaticamente de que existe nova aula no calendário & Aliviada & + Notificação automática, elimina trabalho manual \\ \hline

4 - Conduz sessão ao vivo & Na hora marcada, inicia a transmissão & Confiante & - Pode haver algum problema com a plataforma de reuniões \\ \hline

5 - Modera a conversa & Gerencia quem pode falar, aprova pedidos de palavra dos sêniors & Atenta, levemente tensa & - É necessário ter controle de várias coisas \\ \hline

6 - Finaliza a aula & Após a sessão, acessa relatório automático: presença, histórico por membro & Satisfeita & + Relatório gerado automaticamente \newline\newline + Visão por atividade e por usuário \\ \hline
\end{tabular}
\end{table}